\begin{textsample}{POS Dim 1 – vem\_ed\_21 – Score 6.00 – vem\_ed\_21/t105.txt}  \label{ex:f1_pos_168}
\textbf{continuação} A equipe dos Projetos Colaborativos Internacionais (PCIs / Cesu) apresentou presencialmente duas oficinas, três comunicações orais e fez quatro moderações, resumidas nas páginas a \textbf{seguir}.
Osvaldo Succi Junior, coordenador dos PCIs / Cesu, compôs o painel de encerramento da conferência, “What is the future of COIL / Intercâmbio Virtual”, com Jon Rubin (COIL Virtual Exchange Foundation / EUA), Mirjam Hauck (The Open University / Inglaterra) e a participação remota de Samia Chasi (International Education Association of South Africa / África do Sul).
A moderação foi de Lavern Samuels (Durban University of Technology / África do Sul).
No painel, Succi Junior enfatizou a \textbf{importância} de adicionar metodologias compatíveis à abordagem do COIL para atender melhor os perfis de professores e as \textbf{necessidades} dos alunos.
Ressaltou \textbf{que} o falante nativo precisa \textbf{aprender} a conversar com seus parceiros internacionais (devagar, sem gírias) e \textbf{destacou} a \textbf{importância} de \textbf{conectar} os PCIs com a comunidade (entorno das Fatecs, prefeituras, empresas).

% matched lemmas: aprender, conectar, continuação, destacar, importância, necessidade, que, seguir
\end{textsample}
